% Options for packages loaded elsewhere
\PassOptionsToPackage{unicode}{hyperref}
\PassOptionsToPackage{hyphens}{url}
%
\documentclass[
]{article}
\usepackage{amsmath,amssymb}
\usepackage{iftex}
\ifPDFTeX
  \usepackage[T1]{fontenc}
  \usepackage[utf8]{inputenc}
  \usepackage{textcomp} % provide euro and other symbols
\else % if luatex or xetex
  \usepackage{unicode-math} % this also loads fontspec
  \defaultfontfeatures{Scale=MatchLowercase}
  \defaultfontfeatures[\rmfamily]{Ligatures=TeX,Scale=1}
\fi
\usepackage{lmodern}
\ifPDFTeX\else
  % xetex/luatex font selection
\fi
% Use upquote if available, for straight quotes in verbatim environments
\IfFileExists{upquote.sty}{\usepackage{upquote}}{}
\IfFileExists{microtype.sty}{% use microtype if available
  \usepackage[]{microtype}
  \UseMicrotypeSet[protrusion]{basicmath} % disable protrusion for tt fonts
}{}
\makeatletter
\@ifundefined{KOMAClassName}{% if non-KOMA class
  \IfFileExists{parskip.sty}{%
    \usepackage{parskip}
  }{% else
    \setlength{\parindent}{0pt}
    \setlength{\parskip}{6pt plus 2pt minus 1pt}}
}{% if KOMA class
  \KOMAoptions{parskip=half}}
\makeatother
\usepackage{xcolor}
\usepackage{color}
\usepackage{fancyvrb}
\newcommand{\VerbBar}{|}
\newcommand{\VERB}{\Verb[commandchars=\\\{\}]}
\DefineVerbatimEnvironment{Highlighting}{Verbatim}{commandchars=\\\{\}}
% Add ',fontsize=\small' for more characters per line
\usepackage{framed}
\definecolor{shadecolor}{RGB}{248,248,248}
\newenvironment{Shaded}{\begin{snugshade}}{\end{snugshade}}
\newcommand{\AlertTok}[1]{\textcolor[rgb]{0.94,0.16,0.16}{#1}}
\newcommand{\AnnotationTok}[1]{\textcolor[rgb]{0.56,0.35,0.01}{\textbf{\textit{#1}}}}
\newcommand{\AttributeTok}[1]{\textcolor[rgb]{0.13,0.29,0.53}{#1}}
\newcommand{\BaseNTok}[1]{\textcolor[rgb]{0.00,0.00,0.81}{#1}}
\newcommand{\BuiltInTok}[1]{#1}
\newcommand{\CharTok}[1]{\textcolor[rgb]{0.31,0.60,0.02}{#1}}
\newcommand{\CommentTok}[1]{\textcolor[rgb]{0.56,0.35,0.01}{\textit{#1}}}
\newcommand{\CommentVarTok}[1]{\textcolor[rgb]{0.56,0.35,0.01}{\textbf{\textit{#1}}}}
\newcommand{\ConstantTok}[1]{\textcolor[rgb]{0.56,0.35,0.01}{#1}}
\newcommand{\ControlFlowTok}[1]{\textcolor[rgb]{0.13,0.29,0.53}{\textbf{#1}}}
\newcommand{\DataTypeTok}[1]{\textcolor[rgb]{0.13,0.29,0.53}{#1}}
\newcommand{\DecValTok}[1]{\textcolor[rgb]{0.00,0.00,0.81}{#1}}
\newcommand{\DocumentationTok}[1]{\textcolor[rgb]{0.56,0.35,0.01}{\textbf{\textit{#1}}}}
\newcommand{\ErrorTok}[1]{\textcolor[rgb]{0.64,0.00,0.00}{\textbf{#1}}}
\newcommand{\ExtensionTok}[1]{#1}
\newcommand{\FloatTok}[1]{\textcolor[rgb]{0.00,0.00,0.81}{#1}}
\newcommand{\FunctionTok}[1]{\textcolor[rgb]{0.13,0.29,0.53}{\textbf{#1}}}
\newcommand{\ImportTok}[1]{#1}
\newcommand{\InformationTok}[1]{\textcolor[rgb]{0.56,0.35,0.01}{\textbf{\textit{#1}}}}
\newcommand{\KeywordTok}[1]{\textcolor[rgb]{0.13,0.29,0.53}{\textbf{#1}}}
\newcommand{\NormalTok}[1]{#1}
\newcommand{\OperatorTok}[1]{\textcolor[rgb]{0.81,0.36,0.00}{\textbf{#1}}}
\newcommand{\OtherTok}[1]{\textcolor[rgb]{0.56,0.35,0.01}{#1}}
\newcommand{\PreprocessorTok}[1]{\textcolor[rgb]{0.56,0.35,0.01}{\textit{#1}}}
\newcommand{\RegionMarkerTok}[1]{#1}
\newcommand{\SpecialCharTok}[1]{\textcolor[rgb]{0.81,0.36,0.00}{\textbf{#1}}}
\newcommand{\SpecialStringTok}[1]{\textcolor[rgb]{0.31,0.60,0.02}{#1}}
\newcommand{\StringTok}[1]{\textcolor[rgb]{0.31,0.60,0.02}{#1}}
\newcommand{\VariableTok}[1]{\textcolor[rgb]{0.00,0.00,0.00}{#1}}
\newcommand{\VerbatimStringTok}[1]{\textcolor[rgb]{0.31,0.60,0.02}{#1}}
\newcommand{\WarningTok}[1]{\textcolor[rgb]{0.56,0.35,0.01}{\textbf{\textit{#1}}}}
\usepackage{longtable,booktabs,array}
\usepackage{calc} % for calculating minipage widths
% Correct order of tables after \paragraph or \subparagraph
\usepackage{etoolbox}
\makeatletter
\patchcmd\longtable{\par}{\if@noskipsec\mbox{}\fi\par}{}{}
\makeatother
% Allow footnotes in longtable head/foot
\IfFileExists{footnotehyper.sty}{\usepackage{footnotehyper}}{\usepackage{footnote}}
\makesavenoteenv{longtable}
\usepackage{graphicx}
\makeatletter
\def\maxwidth{\ifdim\Gin@nat@width>\linewidth\linewidth\else\Gin@nat@width\fi}
\def\maxheight{\ifdim\Gin@nat@height>\textheight\textheight\else\Gin@nat@height\fi}
\makeatother
% Scale images if necessary, so that they will not overflow the page
% margins by default, and it is still possible to overwrite the defaults
% using explicit options in \includegraphics[width, height, ...]{}
\setkeys{Gin}{width=\maxwidth,height=\maxheight,keepaspectratio}
% Set default figure placement to htbp
\makeatletter
\def\fps@figure{htbp}
\makeatother
\setlength{\emergencystretch}{3em} % prevent overfull lines
\providecommand{\tightlist}{%
  \setlength{\itemsep}{0pt}\setlength{\parskip}{0pt}}
\setcounter{secnumdepth}{-\maxdimen} % remove section numbering
\usepackage{fancyhdr}
\usepackage{xurl}
\usepackage{fvextra}
\usepackage{microtype}
\usepackage{titlesec}
\definecolor{AzulProjeto}{HTML}{164B6B}
\definecolor{CinzaProjeto}{HTML}{586A76}
\titleformat{\section}{\Large\bfseries\color{AzulProjeto}}{\thesection}{0.6em}{}
\titleformat{\subsection}{\large\bfseries\color{AzulProjeto}}{\thesubsection}{0.6em}{}
\titleformat{\subsubsection}{\normalsize\bfseries\color{AzulProjeto}}{\thesubsubsection}{0.6em}{}
\pagestyle{fancy}
\fancyhf{}
\fancyhead[L]{\small\color{CinzaProjeto}INDUTORES ACOPLADOS | CONVERSOR DE 2 kW}
\fancyfoot[L]{\footnotesize\color{CinzaProjeto}Cálculos e hipóteses sujeitos à validação experimental}
\fancyfoot[R]{\small\thepage}
\setlength{\headheight}{18pt}
\setlength{\emergencystretch}{4em}
\setlength{\tabcolsep}{4pt}
\renewcommand{\arraystretch}{1.15}
\AtBeginEnvironment{longtable}{\footnotesize}
\RecustomVerbatimEnvironment{Verbatim}{Verbatim}{breaklines=true,breakanywhere=true,fontsize=\small}
\AtBeginDocument{\hypersetup{colorlinks=true,linkcolor=AzulProjeto,urlcolor=AzulProjeto}}
\ifLuaTeX
  \usepackage{selnolig}  % disable illegal ligatures
\fi
\usepackage{bookmark}
\IfFileExists{xurl.sty}{\usepackage{xurl}}{} % add URL line breaks if available
\urlstyle{same}
\hypersetup{
  hidelinks,
  pdfcreator={LaTeX via pandoc}}

\author{}
\date{}

\begin{document}

\subsection{1. Finalidade e limites do
exemplo}\label{finalidade-e-limites-do-exemplo}

Este documento resolve um exemplo de dois indutores acoplados idênticos
para estudar o conversor de 2 kW, com Lm=209 µH e relação Ns/Np=2/3. Os
cenários A (50 kHz) e B (100 kHz) são comparados sem transformar
hipóteses em resultados de simulação.

\textbf{O exemplo usa uma geometria hipotética maior. Não corresponde a
um E55, PQ40 ou outro núcleo de estoque confirmado. Não é uma liberação
para construir e energizar os componentes nominais do conversor.} Ele
mostra como fechar uma ficha numérica, reprovar tentativas, ajustar
espiras/janela e definir os ensaios que faltam. A seleção comercial e a
operação em potência exigem os dados ausentes.

O objetivo didático é avançar além de uma lista de fórmulas: escolher
valores explícitos, executar cada substituição e montar uma sequência
física coerente, mantendo a distinção entre dados medidos, requisitos e
escolhas de estudo.

\subsection{2. O que é dado e o que foi
adotado}\label{o-que-uxe9-dado-e-o-que-foi-adotado}

\begin{longtable}[]{@{}
  >{\raggedright\arraybackslash}p{(\columnwidth - 4\tabcolsep) * \real{0.3000}}
  >{\raggedleft\arraybackslash}p{(\columnwidth - 4\tabcolsep) * \real{0.4000}}
  >{\raggedright\arraybackslash}p{(\columnwidth - 4\tabcolsep) * \real{0.3000}}@{}}
\toprule\noalign{}
\begin{minipage}[b]{\linewidth}\raggedright
Parâmetro
\end{minipage} & \begin{minipage}[b]{\linewidth}\raggedleft
Valor
\end{minipage} & \begin{minipage}[b]{\linewidth}\raggedright
Situação
\end{minipage} \\
\midrule\noalign{}
\endhead
\bottomrule\noalign{}
\endlastfoot
Potência nominal total & 2.000 W & Requisito do conversor; não
multiplicar pelas saídas \\
Entrada & 56 Vcc & Requisito \\
Barramento & 340 V / +170 V / −170 V & Conexões e cargas ainda precisam
de esquema completo \\
Lm por componente & 209 µH & Requisito; referência primária ainda a
confirmar no PSIM \\
Ns/Np & 2/3 & Aproxima 0,667 \\
Ipk de referência & 7,83 A & Arredondamento do maior pico do arquivo \\
Envelope de estudo & 10 A & Hipótese de projeto, não transiente
simulado \\
Ip,rms / Is,rms físicos & 8 A / 6 A & Hipóteses didáticas; não extraídas
de iM \\
Ae / Aw úteis & 354 mm² / 400 mm² & Geometria fictícia; não dado de
catálogo \\
le / μr & 120 mm / 2.000 & Hipóteses para relutância \\
Janela útil & 25 mm × 16 mm & Gabarito hipotético, com convenção de uma
passagem das espiras \\
Blim & 0,20 T & Meta de estudo; não Bsat do material \\
Jmáx / ku de projeto & 3 A/mm² / 0,35 & Critérios preliminares \\
Tolerâncias de fluxo & L +5\%; Ae −5\% & Hipóteses de robustez a
verificar no catálogo \\
Temperatura do cobre para R & 100 °C & Condição de cálculo da
resistência, não previsão térmica \\
Material ferrítico & Não especificado & Seleção depende de curvas
oficiais e estoque \\
\end{longtable}

Ae é seção magnética efetiva; Aw é a seção útil de uma janela
atravessada por todos os condutores do desenho. Se um fabricante
publicar a soma de duas janelas, converter antes de comparar. Não contar
o mesmo cobre duas vezes nem duplicar área útil.

\subsection{3. Confirmar o circuito antes do
bobinamento}\label{confirmar-o-circuito-antes-do-bobinamento}

O recorte recebido sugere dois componentes independentes, cada um com P
e S, ramo Lm e dispersão Lk. A imagem parcial não confirma o modelo
completo ou os sentidos dos pontos. Para o exemplo, cada conjunto é
construído separadamente e usa a mesma ficha; a ligação no conversor
deve seguir o esquema completo.

Com correntes definidas entrando nos pontos, adotar iM=iP+n·iS. Portanto
iP=iM−n·iS. A corrente eficaz física exige integrar iP(t) e iS(t),
incluindo seus intervalos de condução. Não usar 6,18 A de iM como se
fosse automaticamente a corrente do primário.

Os valores 8/6 A deste documento apenas permitem demonstrar o cálculo de
cobre. Não são garantia de atendimento ao conversor de 2 kW. Uma vez
exportadas as correntes físicas, substituí-las no código; a contagem de
fios e a janela podem mudar.

\subsection{4. Conferência das correntes
fornecidas}\label{conferuxeancia-das-correntes-fornecidas}

Janela de cinco períodos: 34,940 a 35,040 ms. O arquivo tem 1.504
amostras e o período nominal é 20 µs. O programa integra no tempo,
interpola os limites e calcula o RMS pelo quadrado de cada trecho
linear.

\begin{longtable}[]{@{}lrr@{}}
\toprule\noalign{}
Grandeza & Lm1 & Lm2 \\
\midrule\noalign{}
\endhead
\bottomrule\noalign{}
\endlastfoot
Média (A) & 6,099703 & 5,966345 \\
RMS magnetizante (A) & 6,180803 & 6,049461 \\
Máximo na janela (A) & 7,823397 & 7,691064 \\
Mínimo na janela (A) & 4,373278 & 4,237980 \\
Pico a pico (A) & 3,450119 & 3,453084 \\
Ondulação (\%) & 56,5621 & 57,8760 \\
\end{longtable}

Maior valor de todo o arquivo: 7,827687 A. Não é máximo garantido de
partida ou desequilíbrio.

\begin{figure}
\centering
\includegraphics{figuras/correntes_50k.png}
\caption{Correntes magnetizantes medidas na exportação do PSIM.}
\end{figure}

A área de tensão de uma rampa é aproximadamente L·Δi. Para Lm1:
209×10⁻⁶×3,450119=721,075×10⁻⁶ V·s. O valor usa o envelope de cinco
ciclos; exportar vM(t) e analisar cada ciclo para o fechamento
definitivo.

\subsection{5. Por que o exemplo inicial 24/16 foi
revisto}\label{por-que-o-exemplo-inicial-2416-foi-revisto}

Np≥L·Ipk/(Ae·Blim). Com apenas o pico 7,83 A:

\texttt{Np≥209×10⁻⁶×7,83/(354×10⁻⁶×0,20)=23,115}.

24/16 satisfaz a razão, mas B≈0,1926 T e há pouca folga até a meta de
0,20 T. A 10 A, B≈0,2460 T. Essa tentativa não é adequada ao envelope
didático escolhido.

Para 10 A e sem tolerâncias: Np≥29,52; a opção 30/20 chega a 0,1968 T.
Com L+5\% e Ae−5\%, sobe a 0,2175 T e reprova a meta. A opção 33/22
reduz o fluxo, mas não permite duas metades iguais de P. Para facilitar
intercalamento simétrico e manter folga, escolher \textbf{Np=36 e
Ns=24}.

Na janela antiga hipotética de 250 mm², 36/24 com os mesmos feixes não
cabe no critério de ocupação usado. O exemplo aumenta Aw para 400 mm².
Isso é uma decisão geométrica de estudo, não uma alteração presumida de
um núcleo comercial.

\subsection{6. Seleção por AeAw}\label{seleuxe7uxe3o-por-aeaw}

O cobre exige:

\texttt{Aw·ku≥Np·Ip,rms/J+Ns·Is,rms/J}.

Como Ns=n·Np e Np≥L·Ipk/(Ae·Blim):

\texttt{Ae·Aw≥L·Ipk·(Ip,rms+n·Is,rms)/(Blim·J·ku)}.

Substituição SI:

\texttt{AeAw\_req\ =\ (209×10⁻⁶×10×(8+(2/3)×6))/(0,20×3×10⁶×0,35)}.

Resultado: \textbf{11,9429 cm⁴}.

Área disponível: 354×400=141.600 mm⁴=\textbf{14,16 cm⁴}. A triagem
passa: 14,16\textgreater11,9429. O cálculo de camadas, a margem
magnética e o calor continuam obrigatórios.

O método foi derivado do fluxo e da área de cobre, sem usar a fórmula de
potência de um Flyback. A corrente de pico é magnetizante; as correntes
do somatório são as dos enrolamentos físicos. ku já representa a
utilização global de cobre nu no critério.

\subsection{7. Espiras, fluxo e energia}\label{espiras-fluxo-e-energia}

Escolha inteira: P=36 espiras; S=24; n=24/36=0,666667. Primário dividido
em P1=18 e P2=18 em série aditiva.

\texttt{AL\_alvo=L/Np²=209\ µH/1.296=161,265432\ nH/espira²}.

\texttt{Ls,mag=(24/36)²×209=92,888889\ µH}.

O entreferro será ajustado para esse AL. O arredondamento das espiras
sem ajuste de AL não garante 209 µH.

\begin{longtable}[]{@{}llr@{}}
\toprule\noalign{}
Condição & Cálculo & B \\
\midrule\noalign{}
\endhead
\bottomrule\noalign{}
\endlastfoot
Pico do arquivo arredondado & 209 µH×7,83/(36×354 mm²) & 0,128411 T \\
Envelope hipotético & 209 µH×10/(36×354 mm²) & 0,163999 T \\
Envelope, L+5\%, Ae−5\% & 219,45 µH×10/(36×336,3 mm²) & 0,181262 T \\
\end{longtable}

O último valor é inferior à meta 0,20 T. Ainda falta verificar Bsat
mínimo a quente e fluxo local com Amin da peça real. Ae efetivo não
substitui a menor seção local, e μ não foi comprovado em polarização CC.

Energia no envelope: \texttt{W=½×209×10⁻⁶×10²=10,45\ mJ} por unidade. No
pico de referência, 6,407 mJ. Não somar novamente a energia dos dois
enrolamentos como se fossem dois núcleos independentes: ambos
compartilham a energia do mesmo componente.

A energia do modelo de dois indutores acoplados pode ser escrita
\texttt{½Lp\ ip²+½Ls\ is²+M\ ip\ is}, respeitando sentidos. O modelo Lm
mais dispersões distribui essa energia em ramo magnetizante e fluxos de
dispersão.

\subsection{8. Comparação magnética entre
frequências}\label{comparauxe7uxe3o-magnuxe9tica-entre-frequuxeancias}

A 50 kHz, \texttt{ΔB=L·Δi/(Np·Ae)}:

\texttt{ΔB1=209×10⁻⁶×3,450119/(36×354×10⁻⁶)=56,5815\ mT}.

B é unipolar com forte polarização CC; a amplitude AC aproximadamente
metade de ΔB não é igual a Bpk total.

Para 100 kHz, apenas se tensões por estado, duty e médias se mantiverem,
Δi e ΔB cairiam à metade. O projeto B mantém 36/24 e 209 µH neste
exemplo comparativo; não se divide N por dois.

\begin{longtable}[]{@{}
  >{\raggedright\arraybackslash}p{(\columnwidth - 4\tabcolsep) * \real{0.2727}}
  >{\raggedleft\arraybackslash}p{(\columnwidth - 4\tabcolsep) * \real{0.3636}}
  >{\raggedleft\arraybackslash}p{(\columnwidth - 4\tabcolsep) * \real{0.3636}}@{}}
\toprule\noalign{}
\begin{minipage}[b]{\linewidth}\raggedright
Grandeza magnetizante
\end{minipage} & \begin{minipage}[b]{\linewidth}\raggedleft
A --- 50 kHz
\end{minipage} & \begin{minipage}[b]{\linewidth}\raggedleft
B --- 100 kHz, condicional
\end{minipage} \\
\midrule\noalign{}
\endhead
\bottomrule\noalign{}
\endlastfoot
Δi1 (A) & 3,450119 & 1,725060 \\
Δi2 (A) & 3,453084 & 1,726542 \\
Ondulação 1 (\%) & 56,5621 & 28,2810 \\
Ondulação 2 (\%) & 57,8760 & 28,9380 \\
ΔB1 (mT) & 56,5815 & 28,2908 \\
ΔB2 (mT) & 56,6301 & 28,3151 \\
Irms magnetizante 1 (A) & 6,180803 do arquivo & 6,119997 pela
aproximação triangular \\
Irms magnetizante 2 (A) & 6,049461 do arquivo & 5,987126 pela
aproximação triangular \\
\end{longtable}

O programa também calcula uma aproximação triangular a 50 kHz,
ligeiramente diferente do RMS integrado. Essa diferença não é erro de
integração: são métodos/dados distintos. Para a conferência do PSIM usar
os valores integrados.

\textbf{209 µH continua sem atender 20\%}. A 50 kHz, o cálculo estático
daria 591,1/604,8 µH; a 100 kHz condicional, 295,5/302,4 µH. Não
incorporar essas alterações ao exemplo sem refazer energia, controle,
fluxo, espiras e janela.

\subsection{9. Entreferro calculado e significado
físico}\label{entreferro-calculado-e-significado-fuxedsico}

\texttt{Rtotal=Np²/Lm=1.296/(209×10⁻⁶)=6,20096×10⁶\ A/Wb}.

\texttt{g\_eq=μ0\ Ae\ Np²/Lm−le/μr}.

\texttt{g\_eq=(4π×10⁻⁷×354×10⁻⁶×1.296)/(209×10⁻⁶)−0,120/2.000}.

Resultado: \textbf{2,698493 mm} sem espraiamento.

A relutância do núcleo foi incluída, mas le/μr e a geometria são
hipotéticos. Com gap apenas na coluna central e faces externas
encostadas, esse é aproximadamente o comprimento central somado. Com
espaçamento uniforme de toda a junta e áreas externas totais
equivalentes à central, o fluxo cruza duas juntas em série; a espessura
inicial de espaçador seria aproximadamente \textbf{1,34925 mm}, não
2,69849 mm em cada face.

Esse gap é grande o suficiente para tornar o fringing uma pendência
relevante. Não se propõe simplesmente colocá-lo ao lado do cobre e
declarar perdas baixas. Para uma peça real, procurar gap
central/distribuído apropriado do fabricante, prever distância do cobre
e avaliar aquecimento local; o tamanho físico pode mudar ao incluir Agap
efetiva. Se o fringing impedir a margem térmica, a construção proposta
deve ser rejeitada ou reformulada.

AL é ajustado medindo com S aberto e pressão de montagem reproduzível,
depois verificando a relação entre indutância própria e Lm do modelo.
Pequeno sinal em i≈0 não comprova a indutância incremental entre 4 e 10
A.

\subsection{10. Cobre: área e quantidade de
filamentos}\label{cobre-uxe1rea-e-quantidade-de-filamentos}

Com J≤3 A/mm²: AcuP≥8/3=2,6667 mm² e AcuS≥6/3=2,0000 mm².

AWG30: \texttt{d=0,127×92\^{}((36−30)/39)=0,254639\ mm}.

\texttt{Afio=πd²/4=0,050926\ mm²}.

Número mínimo de fios: \texttt{mP=ceil(2,6667/0,050926)=53};
\texttt{mS=ceil(2/0,050926)=40}.

Áreas finais P=2,699079 mm², S=2,037041 mm². Densidades P=2,9640 A/mm² e
S=2,9454 A/mm².

Fios esmaltados em paralelo precisam ter comprimentos e percursos
semelhantes, transposição das posições e conexão de todos os filamentos.
Litz exige construção e terminação apropriadas. Muitos fios não garantem
Rac baixo se não houver controle de proximidade e divisão de corrente.

A profundidade do cobre a 20 °C é cerca de 0,296 mm a 50 kHz e 0,209 mm
a 100 kHz. O raio do AWG30 é 0,127 mm, inferior a ambas. Isso ajuda no
efeito pelicular individual, mas as três camadas de cada seção ainda
estão sujeitas a proximidade.

Comparar AWG28 (33/25 fios para as mesmas hipóteses), AWG30 (53/40) e
AWG32 (84/63). AWG32 pode reduzir perdas de campo em certas condições,
mas aumenta esmalte e dificuldade de fabricação. Testar feixe real antes
de comprar toda a metragem.

\subsection{11. Janela, camadas e desenho do
carretel}\label{janela-camadas-e-desenho-do-carretel}

Adotar diâmetro externo máximo hipotético de 0,30 mm. Feixe P: 53 fios
organizados em retângulo aproximado de 3,3×1,5 mm (55 posições, duas
vazias). Feixe S: 40 fios, 3,0×1,2 mm. Essas dimensões são de gabarito
ordenado, não de feixe circular solto.

Largura axial útil=25 mm. Então P comporta floor(25/3,3)=7
espiras/camada; S comporta floor(25/3,0)=8.

\begin{longtable}[]{@{}lrrl@{}}
\toprule\noalign{}
Seção radial & Espiras & Fios paralelos & Distribuição por camada \\
\midrule\noalign{}
\endhead
\bottomrule\noalign{}
\endlastfoot
P1, interna & 18 & 53×AWG30 & 7 + 7 + 4 \\
S, intermediária & 24 & 40×AWG30 & 8 + 8 + 8 \\
P2, externa & 18 & 53×AWG30 & 7 + 7 + 4 \\
\end{longtable}

Fios: 3×1,5+3×1,2+3×1,5=12,6 mm radiais. Reservas didáticas: 0,20 mm
inicial; duas barreiras P/S de 0,50 mm; seis fitas entre camadas de 0,05
mm; 0,20 mm externo. Total=1,70 mm. Altura final=\textbf{14,30 mm};
folga até 16 mm=\textbf{1,70 mm}.

\textbf{Essas espessuras são hipóteses geométricas, não uma
especificação de isolação aprovada.} A largura de 25 mm já deve excluir
margens elétricas e mecânicas. O recorte não determina se existe
isolação de segurança requerida entre P e S. As distâncias e a classe do
sistema devem ser definidas pela aplicação, picos e norma pertinente.

\begin{figure}
\centering
\includegraphics{figuras/arranjo_exemplo.png}
\caption{Corte radial do exemplo, com uma janela útil hipotética.}
\end{figure}

P1 e P2 ficam em série aditiva. Marcar todos os inícios e fins e
confirmar com LCR ou gerador isolado. A ligação dos pontos em Lm1/Lm2 no
conversor pode diferir; o desenho abaixo é de um componente isolado, não
o esquema completo do conversor.

\begin{figure}
\centering
\includegraphics{figuras/terminais_exemplo.png}
\caption{Identificação didática dos terminais e da emenda do primário.}
\end{figure}

\subsection{12. Fator de execução e verificação de
área}\label{fator-de-execuuxe7uxe3o-e-verificauxe7uxe3o-de-uxe1rea}

A definição da aula é FE=Aw\_min/Aw, incluindo o esmalte:

\texttt{Aw\_min=(36×53+24×40)×π×0,30²/4=202,727\ mm²}.

\texttt{FE\_anexo=202,727/400=0,506817}.

Para empacotamento 0,785 e reserva de 42,5 mm² equivalente às
fitas/barreiras da hipótese:

\texttt{FE\_geom=(202,727/0,785+42,5)/400=0,751877}.

Cobre nu real na janela:
\texttt{(36×2,699079+24×2,037041)/400=0,365140}. É maior que o ku de
triagem 0,35, mas inferior à capacidade geométrica assumida com
isolantes. Além da área, a contagem de camadas demonstrou folga radial.

As embalagens retangulares dos feixes ocupam aproximadamente 264,6 mm²,
ligeiramente mais que o cálculo circular corrigido. Somando as reservas
de 42,5 mm² resulta 307,1/400=0,76775. O retângulo envolvente de toda a
altura ocupa 25×14,3/400=0,89375, incluindo vazios nas camadas
incompletas. São métricas diferentes; não misturar seus valores.

A aprovação é apenas do gabarito ideal. Se o feixe real aumentar de
largura/altura, as camadas podem mudar de forma descontínua. Medir uma
amostra transposta e rever sem apertar o esmalte.

\subsection{13. Comprimentos por camada e
resistência}\label{comprimentos-por-camada-e-resistuxeancia}

Usar o mesmo MLT para a seção interna e externa subestima a resistência
externa. O exemplo utiliza contorno convexo hipotético com perímetro
base C0=114 mm. A linha paralela a uma distância radial r tem perímetro
aproximado:

\texttt{C(r)=114+2πr\ mm}.

O contorno é um modelo didático; na peça real o perímetro deve ser
medido ou obtido do desenho. Os centros das camadas incluem fita
inicial, barreiras e fitas entre camadas.

\begin{longtable}[]{@{}llr@{}}
\toprule\noalign{}
Seção & Raios das camadas (mm) & MLT médio ponderado (mm) \\
\midrule\noalign{}
\endhead
\bottomrule\noalign{}
\endlastfoot
P1 & 0,95; 2,50; 4,05 & 128,084807 \\
S & 5,90; 7,15; 8,40 & 158,924775 \\
P2 & 10,25; 11,80; 13,35 & 186,518430 \\
\end{longtable}

Por exemplo, MLT\_P1={[}7C(0,95)+7C(2,50)+4C(4,05){]}/18. Comprimento de
P inclui as duas seções e 0,20 m de terminais por filamento; S inclui 24
espiras e 0,20 m.

Resultados: lP=\textbf{5,862858 m}, lS=\textbf{4,014195 m} por
filamento. A emenda P1/P2 pode exigir reserva adicional real.

A 100 °C:
\texttt{ρ=0,01724×{[}1+0,00393×(100−20){]}=0,022660256\ Ω·mm²/m}.

\texttt{RP=ρlP/AcuP=49,221926\ mΩ}.

\texttt{RS=ρlS/AcuS=44,654323\ mΩ}.

\texttt{Pcu,DC=8²RP+6²RS=4,757759\ W} por componente, sob as hipóteses
de corrente e temperatura. Esse é o valor que seria dissipado pelas
mesmas correntes RMS se a resistência fosse Rdc; não é a perda real em
alta frequência.

\subsection{14. Perdas, orçamento térmico e resultado
condicional}\label{perdas-oruxe7amento-tuxe9rmico-e-resultado-condicional}

Sem curvas de ferrite, espectro dos fios ou montagem térmica, não se
determina a perda total real. Para executar um exemplo completo sem
atribuir números falsos a N87/N97, estabelecer \textbf{metas de
aceitação}, que devem ser demonstradas:

\begin{itemize}
\tightlist
\item
  Perda de cobre total ≤1,5×Pcu,DC=7,136638 W por unidade.
\item
  Perda de núcleo ≤3,0 W por unidade.
\item
  Resistência térmica efetiva para o ponto quente ≤5 K/W.
\item
  Ambiente 25 °C e limite de estudo 80 °C.
\end{itemize}

Se todas essas metas forem satisfeitas:

\texttt{Ptotal≤10,136638\ W} por unidade.

\texttt{Tquente≤25+5×10,136638=75,683\ °C}.

O orçamento passaria a meta de 80 °C, mas \textbf{não é previsão medida
ou estimativa oficial do ferrite}. A resistência foi calculada a 100 °C
como hipótese conservadora em relação a esse limite de estudo.
Resistência térmica e temperatura de ponto quente exigem ensaio; a
ferrite e o cobre podem ter temperaturas diferentes.

Para duas unidades: perda combinada ≤20,273 W se os orçamentos forem
demonstrados. A parcela seria cerca de 1,014\% dos 2 kW, sem incluir
semicondutores, capacitores e conexões. Não chamar isso de eficiência do
conversor.

Volume magnético efetivo hipotético Ve≈Ae·le=42,48 cm³. A meta Pcore≤3 W
equivale a Pv≤3/(42,48×10⁻⁶)=70,62 kW/m³, a verificar na forma de onda e
polarização reais. O número é limite permitido, não perda específica do
material.

O cálculo correto do cobre é \texttt{Pcu=Idc²Rdc+ΣIh,rms²Rac(hf,T)}. A
perda de núcleo é \texttt{Pcore=Pv×Ve}, com f, ΔB, forma de onda, Bdc e
temperatura apropriados. Para curvas senoidais, usar amplitude AC, sem
confundir ΔB pico a pico com Bpk da curva. A iGSE e as correções de
polarização só são úteis com parâmetros válidos.

A tabela seguinte é uma sensibilidade SEM associação com 50 ou 100 kHz;
mostra quanto os ensaios precisam controlar as perdas.

\begin{longtable}[]{@{}
  >{\raggedright\arraybackslash}p{(\columnwidth - 8\tabcolsep) * \real{0.1579}}
  >{\raggedleft\arraybackslash}p{(\columnwidth - 8\tabcolsep) * \real{0.2105}}
  >{\raggedleft\arraybackslash}p{(\columnwidth - 8\tabcolsep) * \real{0.2105}}
  >{\raggedleft\arraybackslash}p{(\columnwidth - 8\tabcolsep) * \real{0.2105}}
  >{\raggedleft\arraybackslash}p{(\columnwidth - 8\tabcolsep) * \real{0.2105}}@{}}
\toprule\noalign{}
\begin{minipage}[b]{\linewidth}\raggedright
Fator de cobre sobre referência DC
\end{minipage} & \begin{minipage}[b]{\linewidth}\raggedleft
Pcore adotada (W)
\end{minipage} & \begin{minipage}[b]{\linewidth}\raggedleft
Ptotal (W)
\end{minipage} & \begin{minipage}[b]{\linewidth}\raggedleft
T se Rθ=5 K/W (°C)
\end{minipage} & \begin{minipage}[b]{\linewidth}\raggedleft
T se Rθ=10 K/W (°C)
\end{minipage} \\
\midrule\noalign{}
\endhead
\bottomrule\noalign{}
\endlastfoot
1,3 & 1,0 & 7,185 & 60,93 & 96,85 \\
1,3 & 3,0 & 9,185 & 70,93 & 116,85 \\
1,3 & 5,0 & 11,185 & 80,93 & 136,85 \\
1,5 & 1,0 & 8,137 & 65,68 & 106,37 \\
1,5 & 3,0 & 10,137 & 75,68 & 126,37 \\
1,5 & 5,0 & 12,137 & 85,68 & 146,37 \\
2,0 & 1,0 & 10,516 & 77,58 & 130,16 \\
2,0 & 3,0 & 12,516 & 87,58 & 150,16 \\
2,0 & 5,0 & 14,516 & 97,58 & 170,16 \\
\end{longtable}

Esses cenários não escolhem a melhor frequência. Se a meta térmica
falhar, revisar área de cobre, geometria do gap, distribuição do campo,
janela, material e ventilação. Não compensar apenas elevando o limite de
temperatura sem uma justificativa de sistema.

\subsection{15. Dispersão e
capacitância}\label{dispersuxe3o-e-capacituxe2ncia}

O intercalamento P1--S--P2 foi escolhido para reduzir a separação média
das ampère-espiras. Isso não prova um Lk numérico. Faltam forma do
carretel real, campos do gap e medições; não usar um percentual
arbitrário como resultado calculado.

Metas referidas a P: Lk≤10,45 µH e máximo 20,9 µH. No S: ≤4,644 µH e
máximo 9,289 µH. O secundário tem menos espiras e a referência se
transforma por n², não por n.

No modelo de duas indutâncias próprias, Lsc,P=Lp−M²/Ls=Lp(1−k²). Se
Lp≈Lm, meta de 5\% implica k≥0,97468; para 10\%, k≥0,94868. Confirmar o
modelo antes de aplicar esses números ao PSIM.

Aproximar os enrolamentos aumenta capacitância e corrente de
deslocamento \texttt{i=C\ dv/dt}. As duas barreiras P/S e os múltiplos
planos tornam a resposta de alta frequência um compromisso. Uma
frequência maior pode elevar efeitos de modo comum. Isolação adequada
tem prioridade sobre a tentativa de reduzir Lk.

\subsection{16. Fabricação do exemplo, sem ligação ao conversor
nominal}\label{fabricauxe7uxe3o-do-exemplo-sem-ligauxe7uxe3o-ao-conversor-nominal}

Fazer primeiro uma bobina de prova, desenergizada, usando o gabarito de
25×16 mm. O modelo fictício não pode ser comprado por nome; para uso
real deve ser substituído por um conjunto comercial e verificado.

\begin{enumerate}
\def\labelenumi{\arabic{enumi}.}
\tightlist
\item
  Conferir os desenhos e separar dois conjuntos completos de ferrite e
  carretéis compatíveis, quando identificados. Sem essa seleção, o passo
  físico limita-se ao gabarito de prova.
\item
  Medir largura útil, altura radial, perímetro da primeira camada e
  localização do gap. Confirmar convenção de Aw, margens e distância do
  cobre ao gap.
\item
  Inspecionar trincas e arestas. Preparar fitas, mangas, espaçadores
  calibrados, contador de espiras e folha de registro.
\item
  Confirmar AWG30 pelo diâmetro do cobre e o esmalte máximo. Ensaiar
  feixes com 53 e 40 filamentos; medir suas dimensões transpostas.
\item
  Reservar por filamento aproximadamente 5,863 m para o primário
  completo e 4,014 m para o secundário, antes de margens extras de
  manuseio/emenda. Separar P1 e P2 considerando seus MLT diferentes.
\item
  Preparar todos os filamentos de cada feixe com comprimentos
  semelhantes, identificação e transposição. Todas as pontas precisam
  ser soldadas; contar e inspecionar.
\item
  Aplicar 0,20 mm inicial apenas no gabarito geométrico. Para uma peça
  nominal, usar o sistema de isolação aprovado e recalcular espessuras.
\item
  Enrolar P1: 18 espiras, 53 fios paralelos; primeira camada 7, segunda
  7, terceira 4. Fitas intermediárias conforme a ficha. Contar espiras,
  não filamentos.
\item
  Isolar e encaminhar terminais P1i/P1f sem cruzamentos na região ativa.
  Registrar o sentido visto de uma face de referência fixa.
\item
  Aplicar a barreira P/S do plano de isolação; no modelo de espaço ela
  tem 0,50 mm. A fita não pode suprimir margens de escoamento.
\item
  Enrolar S: 24 espiras de 40 fios, distribuídas 8+8+8. Marcar Si/Sf e
  proteger terminais.
\item
  Aplicar a segunda barreira de 0,50 mm no gabarito ou a espessura real
  homologada.
\item
  Enrolar P2: 18 espiras de 53 fios, 7+7+4. Manter convenção e terminais
  acessíveis para verificar a série aditiva.
\item
  Aplicar isolamento externo, mangas e alívio mecânico. Conferir altura
  final: 14,30 mm no modelo. Se exceder a folga disponível, não
  pressionar para fechar.
\item
  Fazer emenda aditiva P1f--P2i após confirmar polaridade. P tem 36
  espiras totais; não colocar P1/P2 em paralelo.
\item
  Montar o núcleo real com gap preliminar somente após avaliar a
  equivalência e o fringing. Não transferir a lâmina de 1,34925 mm para
  qualquer junta comercial sem conferir a geometria.
\item
  Fixar provisoriamente com pressão uniforme e sem formar uma espira
  condutiva fechada ao redor do fluxo. Não colar antes do ajuste.
\item
  Medir relação, polaridade, indutâncias próprias, dispersão e
  resistência. Ajustar gap com pequenos passos e repetir as medições na
  mesma montagem.
\item
  Após aprovação de pequeno sinal, aplicar fixação definitiva e medir
  novamente. O ensaio sob corrente e o térmico ainda são obrigatórios.
\item
  Repetir o processo para o segundo componente e comparar os resultados.
  Só conectar ao circuito depois de confirmar a ligação de cada ponto no
  esquema completo.
\end{enumerate}

\subsection{17. Plano de ensaios com valores de
referência}\label{plano-de-ensaios-com-valores-de-referuxeancia}

\subsubsection{17.1 LCR e
identificação}\label{lcr-e-identificauxe7uxe3o}

Com componente desconectado e capacitores descarregados: compensar
aberto/curto, registrar frequência/amplitude/modelo série ou paralelo e
temperatura. Medir P com S aberto, depois S com P aberto. Pequena
excitação não deve produzir fluxo elevado.

O modelo ideal sem dispersão daria Lp=209 µH e Ls=92,889 µH. Com
dispersão, são indutâncias próprias diferentes do ramo Lm. Se o PSIM
tiver Lσ em série com o lado medido, não ajustar cegamente a leitura
aberta para 209 µH.

Associar P e S em série aditiva/subtrativa e medir L+ e L−.
\texttt{M=(L+−L−)/4}; \texttt{k=M/√(LpLs)}. Valores ideais: M=139,333
µH; L+=580,556 µH; L−=23,222 µH. As espiras 36/24 mantêm a mesma razão e
os mesmos valores de indutância ideais após ajuste do gap.

Curto-circuitar S com conexão de baixa impedância e medir no P; comparar
a grandeza de dispersão equivalente com a meta de 10,45 µH depois de
mapear o modelo. Repetir no outro lado e referir por n². Usar
frequências do instrumento abaixo da ressonância própria.

Medir DCR por quatro fios. O parâmetro Rs em 100 kHz inclui perdas e não
substitui DCR. Para comparar com este exemplo, corrigir resistência pela
temperatura. Os valores calculados a 100 °C são P=49,222 mΩ e S=44,654
mΩ, sem resistência de emenda extra.

\subsubsection{17.2 Indutância sob corrente e
saturação}\label{indutuxe2ncia-sob-corrente-e-saturauxe7uxe3o}

Com fonte limitada, circuito pulsado com reset e clamp, sonda de
corrente/shunt e osciloscópio diferencial, medir v(t) e i(t). Calcular
Linc≈(v−Ri)/(di/dt) e λ=∫(v−Ri)dt. Aceleração abrupta de di/dt indica
queda de L. Começar com pulsos curtos e baixa energia, garantir
desmagnetização e elevar a corrente em etapas dentro da proteção.

10 A é o envelope didático deste exemplo, não um limite aprovado do
circuito. O teste deve ser definido e supervisionado no laboratório.
Repetir a quente e nos extremos de tolerância. LCR sozinho não executa
essa validação.

\subsubsection{17.3 Isolação e
aquecimento}\label{isolauxe7uxe3o-e-aquecimento}

Tensão e duração de hipot dependem da norma e do sistema de isolação;
não deduzi-las somente de 340 V. Medir isolação com equipamento adequado
e o componente fora da eletrônica.

Operar progressivamente com carga controlada, barreiras e proteção.
Registrar temperatura de núcleo, cobre, terminais e ambiente; acompanhar
equilíbrio dos barramentos e picos de corrente. O objetivo de estudo é
demonstrar Pcu≤7,137 W, Pcore≤3 W e Rθ efetiva≤5 K/W, ou refazer o
orçamento com valores medidos. A leitura externa de um termopar não
representa automaticamente o ponto quente.

Em 340 Vcc há risco letal e energia armazenada; usar desligamento,
descarga e confirmação de ausência de tensão antes de qualquer
intervenção. Ensaios em potência e de isolação exigem pessoal habilitado
e montagem protegida.

\subsection{18. Lista de materiais do exemplo para duas
unidades}\label{lista-de-materiais-do-exemplo-para-duas-unidades}

\begin{longtable}[]{@{}
  >{\raggedright\arraybackslash}p{(\columnwidth - 2\tabcolsep) * \real{0.5000}}
  >{\raggedright\arraybackslash}p{(\columnwidth - 2\tabcolsep) * \real{0.5000}}@{}}
\toprule\noalign{}
\begin{minipage}[b]{\linewidth}\raggedright
Item
\end{minipage} & \begin{minipage}[b]{\linewidth}\raggedright
Quantidade ou condição
\end{minipage} \\
\midrule\noalign{}
\endhead
\bottomrule\noalign{}
\endlastfoot
Núcleos & 2 conjuntos completos (4 metades), código ainda não
selecionado \\
Carretéis & 2, geometria real a confirmar \\
AWG30 para P & 2×53×5,862858=621,463 m de filamento \\
AWG30 para S & 2×40×4,014195=321,136 m de filamento \\
AWG30 total & 942,599 m sem reserva adicional \\
AWG30 com 15\% adicional & 1083,988 m, aproximadamente 1,084 km \\
Cobre nu sem reserva & Aproximadamente 0,430 kg \\
Fixação & 2 conjuntos adequados, sem espira em curto \\
Gap & Jogo de espaçadores/gap de fábrica; valor real ajustado \\
Fitas/mangas & Pelo plano de isolação e área real; estoque a
confirmar \\
Terminais & 8 externos e 2 emendas internas P1/P2 \\
\end{longtable}

A reserva de 15\% é hipótese de planejamento e não substitui a metragem
medida de terminações e transposição. Não comprar núcleos por estes
números sem conferir datasheet e fabricante.

\subsection{19. Comparação final A/B e
decisão}\label{comparauxe7uxe3o-final-ab-e-decisuxe3o}

\begin{longtable}[]{@{}
  >{\raggedright\arraybackslash}p{(\columnwidth - 4\tabcolsep) * \real{0.3333}}
  >{\raggedright\arraybackslash}p{(\columnwidth - 4\tabcolsep) * \real{0.3333}}
  >{\raggedright\arraybackslash}p{(\columnwidth - 4\tabcolsep) * \real{0.3333}}@{}}
\toprule\noalign{}
\begin{minipage}[b]{\linewidth}\raggedright
Parâmetro
\end{minipage} & \begin{minipage}[b]{\linewidth}\raggedright
A --- 50 kHz
\end{minipage} & \begin{minipage}[b]{\linewidth}\raggedright
B --- 100 kHz
\end{minipage} \\
\midrule\noalign{}
\endhead
\bottomrule\noalign{}
\endlastfoot
Geometria didática & Ae354/Aw400 mm² & A mesma \\
Material & A selecionar & A selecionar \\
Np/Ns & 36/24 & 36/24 para comparação \\
Lm & 209 µH & 209 µH \\
Entreferro equivalente sem fringing & 2,69849 mm & O mesmo \\
Condutores do exemplo & P53/S40 AWG30 & Mesmos, somente sob hipóteses de
RMS \\
Envelope e B máximo com tolerâncias & 10 A; 0,181262 T & Mesmo envelope
de estudo \\
ΔB1 & 56,5815 mT & 28,2908 mT condicional \\
FE anexo / geométrico & 0,50682 / 0,75188 & Os mesmos \\
Referência DC de cobre & 4,75776 W a 100 °C e 8/6 A & A mesma se os RMS
forem iguais; hipótese \\
Perda AC/núcleo/temperatura real & A medir/calcular com dados oficiais &
A medir/calcular com nova simulação \\
Dispersão & Meta≤10,45 µH, sem previsão aprovada & Mesma meta,
dependência frequencial a medir \\
Conclusão de eficiência real & Não disponível & Não disponível \\
\end{longtable}

O exemplo fecha os números geométricos e define um orçamento térmico
verificável. Não fecha a seleção comercial, o fringing, a polarização
real ou a eficiência do conversor. O fato de o gabarito caber não
assegura aptidão do núcleo do laboratório.

A decisão prática é validar primeiro o caso de 50 kHz, porque existem
dados numéricos. A comparação a 100 kHz continua necessária; aumento de
frequência pode diminuir excursão/perda de núcleo e aumentar perda AC e
de comutação. A escolha deve usar a mesma potência total, mesmas cargas
e condições de temperatura.

\subsection{20. Reprodução e substituição das
hipóteses}\label{reproduuxe7uxe3o-e-substituiuxe7uxe3o-das-hipuxf3teses}

No diretório do pacote, executar:

\begin{Shaded}
\begin{Highlighting}[]
\ExtensionTok{python} \AttributeTok{{-}m}\NormalTok{ pip install }\AttributeTok{{-}r}\NormalTok{ requirements.txt}
\ExtensionTok{python}\NormalTok{ calculo.py }\AttributeTok{{-}{-}correntes}\NormalTok{ dados/CorrenteIndutor.txt }\AttributeTok{{-}{-}saida}\NormalTok{ resultados}
\ExtensionTok{python}\NormalTok{ exemplo\_completo.py }\AttributeTok{{-}{-}config}\NormalTok{ parametros\_projeto\_exemplo.json }\AttributeTok{{-}{-}saida}\NormalTok{ resultados\_exemplo}
\end{Highlighting}
\end{Shaded}

\texttt{calculo.py} confere o arquivo de 50 kHz e mantém o exemplo
inicial de 24/16 para rastreabilidade. \texttt{exemplo\_completo.py} é o
cálculo deste documento, com 36/24 e geometria ampliada hipotética. O
JSON contém todos os valores adotados e os resultados.

Ao receber dados reais: substituir RMS, envelope, Ae/Aw/le, material,
tolerâncias, diâmetros e dimensões; medir o perímetro por camada;
atualizar plano de isolação; introduzir perdas do material e espectro;
executar validações. Se uma condição de aceitação falhar, repetir o
dimensionamento antes da fabricação nominal.

Referências disponíveis: PROJETO\_INDUTOR.pdf (método, espiras, gap,
condutores e fator de execução), McLyman.pdf (componentes magnéticos,
ocupação e espraiamento) e CorrenteIndutor.txt (dados numéricos de 50
kHz). Os datasheets de fabricantes e o modelo PSIM completo não foram
acessados nesta análise.

\end{document}
